\subsection{Explicit rectangles for finite material orders}
\label{r:sec:material-target-rectangle}

The finite footprint of the material jet gives an explicit answer when
the desired derivative orders are specified. Fix integers \(N,K,m\ge0\)
and put
\[
 J_*=N+K,\qquad q=\max(m,2),\qquad Q=q+K,
 \qquad R=R_{J_*,Q}.
\]
Define
\[
 d_0=\max_{0\le j\le J_*}\|V_j(0)\|_{H^Q},\qquad
 X=\sqrt R,\qquad Z=\sqrt{d_0^2+R}.
\]
The initial recurrence \eqref{r:eq:initial} includes the initial source
derivatives and makes \(d_0\) finite. Each temporal row in this selection
has a continuous \(H^Q\) representative, with
\begin{equation}
\begin{aligned}
 \sup_{0\le t\le T}\|V_j(t)\|_{H^Q}^2
 &\le\|V_j(0)\|_{H^Q}^2
 +2\|V_j\|_{L^2H^{Q+1}}\|V_{j+1}\|_{L^2H^{Q-1}}\\
 &\le d_0^2+R=Z^2.
\end{aligned}
\label{r:eq:material-target-trace}
\end{equation}
This follows by pairing the two adjacent Sobolev coordinates when
differentiating the squared \(H^Q\) norm. The Hilbert triple
\(H^{Q+1}\subset H^Q\subset H^{Q-1}\) also supplies continuity.

Every scalar component of \(D_t^ku\) is a sum of monomials
\[
 c\prod_{\ell=1}^{r}\partial_x^{\beta_\ell}V_{j_\ell},
 \quad 1\le r\le k+1,\quad
 \sum_\ell j_\ell=k+1-r,\quad
 \sum_\ell|\beta_\ell|=r-1.
\]
A temporal derivative increments one temporal index. Transport adds
one velocity factor and one spatial derivative. This proves the
displayed footprint by induction. After \(n\) ordinary temporal
derivatives, its temporal total is \(n+k+1-r\). The unique linear term
is \(V_{n+k}\); every nonlinear term has strictly lower temporal indices.

For integer \(q\ge2\), a scalar Sobolev algebra constant is
\[
 K_q=2^q(2\pi)^{-3/2}
 \left(\pi^{3/2}\frac{\Gamma(q-3/2)}{\Gamma(q)}\right)^{1/2}.
\]
Indeed, the weight inequality
\(\langle\xi\rangle^q\le2^{q-1}
(\langle\eta\rangle^q+\langle\xi-\eta\rangle^q)\), Fourier
convolution, and Cauchy--Schwarz with
\(\int\langle\xi\rangle^{-2q}d\xi
=\pi^{3/2}\Gamma(q-3/2)/\Gamma(q)\) give
\(\|vw\|_{H^q}\le K_q\|v\|_{H^q}\|w\|_{H^q}\).

For each component, the total absolute monomial coefficient in
\(D_t^ku\) is at most \(4^k k!\). A monomial with \(r\) factors
produces at most \(r\) temporal terms and \(3r\) transport terms
at the next material derivative. Applying \(n\) further temporal
derivatives costs at most \((k+1)^n\). Set
\[
 L_{n,k,m}=\sqrt3\,4^k k!(k+1)^n\max(1,K_q)^k.
\]
The factor \(\sqrt3\) combines the three scalar components. Since each
spatial factor has order at most \(K\), its \(H^q\) norm is bounded
by its parent's \(H^Q\) norm. Place one factor in \(L^2_tH^q_x\)
and all others in \(L^\infty_tH^q_x\). For every \(n\le N\),
\(k\le K\), this proves
\begin{equation}
\begin{aligned}
 \|\partial_t^nD_t^ku\|_{L^2(0,T;H^m)}
 &\le L_{n,k,m}X\max(1,Z)^k,\\
 \sup_{0\le t\le T}\|\partial_t^nD_t^ku(t)\|_{H^m}
 &\le L_{n,k,m}Z\max(1,Z)^k.
\end{aligned}
\label{r:eq:material-target-bounds}
\end{equation}
Continuity of the factors and of multiplication in \(H^q\) gives a
continuous \(H^m\) terminal trace. Replacing \(m\) by \(m+b\)
also controls spatial derivatives through order \(b\).
Every desired finite material order thus has a specified finite
rectangle within the full intersection. The constants may depend on
those requested orders.
